\begin{table}[ht!]
\centering
\footnotesize
\caption{The half-year event study estimated separately on workers coded from each register vintage, ages 22--25.}
\label{tab:iv_vintage}
\begin{tabular}{lccc}
\toprule
Period & Coded from 2021 & Coded from 2022 & Coded from 2023 \\
\midrule
2019H1 & $-0.010$ (0.016) & $-0.010$ (0.016) & $-0.005$ (0.015) \\
2019H2 & $-0.037$ (0.017) & $-0.037$ (0.017) & $-0.031$ (0.016) \\
2020H1 & $-0.041$ (0.015) & $-0.041$ (0.015) & $-0.035$ (0.015) \\
2020H2 & $-0.053$ (0.016) & $-0.052$ (0.016) & $-0.047$ (0.015) \\
2021H1 & $-0.021$ (0.011) & $-0.021$ (0.011) & $-0.018$ (0.011) \\
2021H2 & $-0.054$ (0.012) & $-0.054$ (0.012) & $-0.051$ (0.011) \\
2022H1 & $+0.000$ (0.000) & $+0.000$ (0.000) & $+0.000$ (0.000) \\
2022H2 & $-0.047$ (0.009) & $-0.047$ (0.009) & $-0.050$ (0.008) \\
2023H1 & $-0.294$ (0.341) & $-0.201$ (0.120) & $+0.006$ (0.010) \\
2023H2 & $-0.557$ (0.405) & $-0.024$ (0.138) & $-0.018$ (0.012) \\
2024H1 & $-0.344$ (0.463) & $-0.358$ (0.278) & $-0.283$ (0.015) \\
2024H2 & $-0.539$ (0.373) & $-0.643$ (0.268) & $-0.347$ (0.017) \\
2025H1 & $-0.532$ (0.533) & $-1.317$ (0.354) & $-0.587$ (0.026) \\
\bottomrule
\end{tabular}
\begin{minipage}{0.90\textwidth}\footnotesize\vspace{4pt}Each column restricts to workers whose occupation code comes from the stated register year, and runs the submitted employer-level event study on that subsample. \textbf{The columns are not comparable and none of them estimates a treatment effect.} Conditioning on code vintage conditions on how recently a worker's occupation was observed, which is itself a function of tenure, entry and mobility; the subsamples shrink sharply in the later periods, which is why the standard errors explode and the point estimates reach implausible magnitudes. The exercise is reported because the editor asked for it and because the instability is the answer: a design whose estimate depends this strongly on which register year supplied the code is not identified. Source: script 41.\end{minipage}
\end{table}
